\documentclass[11pt]{article}
\usepackage[margin=1in]{geometry}
\usepackage{times}
\usepackage[T1]{fontenc}
\usepackage[utf8]{inputenc}
\usepackage{microtype}
\usepackage{enumitem}
\usepackage{booktabs}
\usepackage{titlesec}
\usepackage[hidelinks]{hyperref}
\titleformat{\section}{\normalfont\large\bfseries}{\thesection.}{0.5em}{}
\titlespacing*{\section}{0pt}{1.1em}{0.35em}
\setlist{nosep,leftmargin=1.5em}

\title{\vspace{-1.6em}\textbf{A Robotic Co-Performer for Four-Hand Piano}\vspace{-0.4em}}
\author{[Your Name]\\ \small Georgia Institute of Technology}
\date{\vspace{-0.7em}\small\today\vspace{-1.2em}}

\begin{document}
\maketitle

\section{Objective}
Build a robot that performs the written \emph{secondo} part of four-hand piano repertoire on a
standard 88-key piano while a human performs the \emph{primo}, with both players at the same
keyboard. The target outcome is a working system that can play through four-hand pieces together
with a human, evaluated by ensemble timing accuracy. Expressive timing negotiation and neural
evaluation are deferred to future work (Sec.~\ref{sec:future}).

\section{Background}
Robotic musicianship concerns machines that collaborate with human musicians rather than merely
generate sound~\cite{weinberg2020robotic}. Piano-playing robots have largely been developed and
evaluated as soloists, where the benchmark is dexterous reproduction of a
score~\cite{zakka2023robopianist}.

The closest prior system is that of Wang, Zhang and Iida, who describe it as ``the first instance of
designing and implementing a collaborative robot expressly engineered to work with humans in the
artistic pursuit of piano playing''~\cite{wang2024cooperative}. A UR5 arm with an anthropomorphic
hand plays chord accompaniment beneath a human-performed melody, using a recurrent network for chord
prediction (93\% accuracy) and a model-predictive controller for temporal alignment. The authors
report the following limitations~\cite{wang2024cooperative}:

\begin{enumerate}[label=\textbf{L\arabic*.}]
  \item ``Due to hardware limitations, we simplified the tuning system'' --- the chord vocabulary is
        seven triads.
  \item ``Due to the hardware limitation (UR5 is not agile/fast enough)'', time allocated to
        keystrokes must be reallocated to lateral chord transitions.
  \item ``Only when all 16 tokens of the melody have been fed into the RNN model can the model give a
        chord prediction \ldots resulting in a lag.''
  \item ``Evaluating whether a listener `enjoys' the music involves psychological and subjective
        analysis'' --- deferred.
\end{enumerate}

The task in~\cite{wang2024cooperative} is improvised homophonic accompaniment: generating chords
beneath an unknown melody. Performing a written second part in ensemble is a different task, and to
our knowledge no robotic system has been reported for it.

\section{Why This Is Not Straightforward}
Ensemble coordination operates within tolerances of tens of milliseconds. Reported detection
thresholds for asynchrony in singing ensembles lie between approximately 10 and
38\,ms~\cite{singingsync2019}. Professional string quartets performing at 157\,bpm show note-onset
asynchrony standard deviations of 24 and 28\,ms~\cite{quartetsync2014}. Studies of networked piano
duets test acoustic transmission latencies of 10, 20 and 40\,ms as the musically relevant
range~\cite{nmpduet2021}.

A physical co-performer accumulates latency at every stage of its signal chain. Published figures
exist only for the perceptual tolerances above; the per-stage figures below are \emph{engineering
estimates to be measured in M1--M2}, not established results:

\begin{center}
\begin{tabular}{lr}
\toprule
Stage & Estimate (to be measured) \\
\midrule
Key press $\rightarrow$ MIDI received & 1--5\,ms \\
Score-follower position update & 10--50\,ms \\
Decision and scheduling & 1--5\,ms \\
Solenoid travel to key contact & 10--30\,ms \\
Key $\rightarrow$ sound (acoustic action) & 5--20\,ms \\
\midrule
Total & $\approx$30--110\,ms \\
\bottomrule
\end{tabular}
\end{center}

\noindent
If these estimates hold, the accumulated latency exceeds the tolerances
in~\cite{singingsync2019,quartetsync2014,nmpduet2021}. Establishing whether this is so, and by how
much, is the first objective of the project (M1).

\section{System}
\textbf{A. Fixed solenoid array.} A fixed actuator per key over the secondo register, plus one
sustain-pedal actuator, on a standard 88-key instrument. (Array extent is a scope parameter:
approximately 44 keys covers the secondo range; full 88-key coverage would additionally permit role
reversal.) A fixed array has no lateral travel, which removes the transition-time constraint L2, and
places no moving structure near the human's hands. Solenoid actuation has been reported as less
expressive than brushless direct-drive actuation for robotic percussion~\cite{nime2020bldc}; that
comparison concerns free-trajectory striking, whereas a fixed array is adopted here because a written
secondo part requires many keys to be held simultaneously. The expressive cost will be measured
rather than assumed.

\textbf{B. Score following.} Real-time alignment to the human's position in the score. This is a
mature area~\cite{raphael2010musicplusone,cont2008antescofo}, and expressive accompaniment systems
that adapt to a soloist's tempo, dynamics and articulation already
exist~\cite{cancinochacon2023accompanion}. \textbf{No algorithmic novelty in score following is
claimed}; the work here is integration under the constraints of a shared physical instrument.
Symbolic score handling builds on existing open tooling~\cite{partitura2022}.

\textbf{C. Shared-keyboard note separation.} Both players' notes arrive on one instrument. Two
candidate methods --- keyboard split to independent MIDI channels, and subtraction of the robot's own
issued commands --- will be implemented and compared.

\section{Milestones}
\begin{enumerate}[label=\textbf{M\arabic*.}]
  \item \textbf{Software only.} Score follower drives the secondo through the digital piano's own
        sound engine, removing actuation latency. Measures the latency budget and establishes whether
        ensemble tolerance is reachable in the most favourable case. A negative result here is
        decisive and obtained at no hardware cost.
  \item \textbf{Actuator array.} 8 $\rightarrow$ 24 $\rightarrow$ full secondo range, plus pedal,
        with a custom driver board. Per-stage latency measured on hardware.
  \item \textbf{Ensemble performance.} Play complete four-hand pieces with a human partner.
\end{enumerate}

\noindent
\textbf{Repertoire.} Tchaikovsky's four-hand arrangements of Russian folk songs provide short, simple
test material; Schubert's \emph{Marches militaires} D.733 and Brahms's \emph{Hungarian Dances} serve
as demonstration pieces.

\section{Evaluation}
The primary measure is note-onset asynchrony between the human and robot parts, compared against
human--human four-hand performance of the same passages recorded under the same conditions. The
tolerances reported in~\cite{singingsync2019,quartetsync2014,nmpduet2021} provide the reference
range. A secondary measure is a structured report from the pianist co-performer; the assessment
question deferred in~\cite{wang2024cooperative} is addressed here from the co-performer's side rather
than the listener's.

\section{Future Work}
\label{sec:future}
Two extensions are deliberately excluded from the present scope. \textbf{(i) Expressive timing.}
Existing accompaniment systems assume a fixed leader, whereas coordination in human piano duets
varies with texture and can involve alternating control~\cite{predictivecoding2026}; modelling this
requires a working system first. \textbf{(ii) Neural evaluation.} EEG hyperscanning against the
published human--human baseline for duetting pianists~\cite{hyperscanning2021} becomes possible once
ensemble performance is achieved.

\bibliographystyle{unsrt}
\bibliography{references}
\end{document}
